\documentclass[12pt, titlepage, a4paper]{article}
\usepackage[a4paper, inner=2.5cm, outer=2.5cm, top=2.25cm, bottom=2.25cm]{geometry} %Lepša postavitev strani
%matematika
\usepackage{amsmath}
\usepackage{amsthm}
\usepackage{amssymb}
\usepackage{empheq}
%
\usepackage[slovene]{babel}
\usepackage{hyperref}
\usepackage{nameref}

\usepackage{bussproofs}
\usepackage{graphicx} 
\usepackage{subcaption}

%

% Okrajšave, ker je v dokumentaciji orajšano, in je lažje primerjati kodo z dokumentacijo
\newcommand{\AXC}[1]{\AxiomC{#1}}
\newcommand{\UIC}[1]{\UnaryInfC{#1}}
\newcommand{\BIC}[1]{\BinaryInfC{#1}}
\newcommand{\TIC}[1]{\TrinaryInfC{#1}}
\newcommand{\RL}[1]{\RightLabel{#1}}
\newcommand{\DP}{\DisplayProof}


\newcommand{\pravilo}[1]{\textbf{#1.}}

\newcommand{\G}{\Gamma}
\newcommand{\D}{\Delta}
\newcommand{\la}{\lambda}
\newcommand{\cd}[1]{\texttt{#1}}
\newcommand{\type}{\text{ type}}
\newcommand{\rampa}{\(\vdash\) }

% \newcommand{\darule}{\noindent\hdashrule{\textwidth}{0.4pt}}
\newcommand{\darule}{\hdashrule[0.5ex][x]{\linewidth}{0.5pt}{1.5mm}}

\newenvironment{bprooftree}
  {\leavevmode\hbox\bgroup}
  {\DisplayProof\egroup}



\begin{document}

\section{Linearni sistemi tipov}

\begin{itemize}
    \item Najprej logika
    \item Logika in množice
    \item pravila sklepanja
    \item Modeliranje sveta z logiko
\end{itemize}


\section{Naravna dedukcija}

\begin{itemize}
    \item Tej notaiji se reče naravna dedukcija
    \item Razlaga notacije (premisa, skelp, kontekst, rampa, sodba)
\\
\begin{prooftree}
    \AXC{\(\G \vdash A\)}
    \AXC{\(\G \vdash B\)}
    \RL{ime}
    \BIC{\(\G \vdash A \land B\)}
\end{prooftree}
\item Naravna dedukcija(pravila za konjunkcijo, disjunkcijo, implikacijo)
\\
\\
% \item Aksiom
% \begin{prooftree}
%     \AXC{}
%     \RL{\emph{Ax}}
%     \UIC{\(\G, A \vdash A\)}
% \end{prooftree}
\item   konjunkcija
\begin{prooftree}
    \AXC{\(\G \vdash A\)}
    \AXC{\(\G \vdash B\)}
    \RL{\(\land I\)}
    \BIC{\(\G \vdash A \land B\)}
\end{prooftree}
\begin{equation*}
    \begin{bprooftree}
        \AXC{\(\G \vdash A \land B\)}
        \RL{\(\land E_L\)}
        \UIC{\(\G \vdash A\)}
    \end{bprooftree}
    \qquad
     \begin{bprooftree}
        \AXC{\(\G \vdash A \land B\)}
        \RL{\(\land E_D\)}
        \UIC{\(\G \vdash A\)}
    \end{bprooftree}
\end{equation*}

\item Implikacija
\begin{prooftree}
    \AXC{\(\G, A \vdash B\)}
    \RL{\(\Rightarrow I\)}
    \UIC{\(\G \vdash A \implies B\)}
\end{prooftree}

\begin{prooftree}
    \AXC{\(\G \vdash A \implies B\)}
    \AXC{\(\G \vdash A\)}
    \RL{\(\Rightarrow E\)}
    \BIC{\(\G \vdash B\)}
\end{prooftree}

\item \(A \land B \Rightarrow B \land A\)


\begin{prooftree}
    \AXC{}
    \RL{\(Ax\)}
    \UIC{\(A \land B \vdash A \land B\)}
    \RL{\(\land E_2\)}
    \UIC{\(A \land B \vdash B\)}
    \AXC{}
    \RL{\(Ax\)}
    \UIC{\(A \land B \vdash A \land B\)}
    \RL{\(\land E_1\)}
    \UIC{\(A \land B \vdash A\)}
    \RL{\(\land I\)}
    \BIC{\(A \land B \vdash B \land A\)}
    \RL{\(\Rightarrow I\)}
    \UIC{\(\vdash A \land B \Rightarrow B \land A\)}
\end{prooftree}


% \item disjunkcija
% \begin{equation*}
%     \begin{bprooftree}
%         \AXC{\(\G \vdash A\)}
%         \RL{\(\lor I_L\)}
%         \UIC{\(\G \vdash A \lor B\)}
%     \end{bprooftree}
%     \begin{bprooftree}
%         \AXC{\(\G \vdash B\)}
%         \RL{\(\lor I_D\)}
%         \UIC{\(\G \vdash A \lor B\)}
%     \end{bprooftree}
% \end{equation*}
% %
% %
% \begin{prooftree}
%     \AXC{\(\G \vdash A \lor B\)}
% %
%     \AXC{\(\G, A \vdash C\)}
% %
%     \AXC{\(\G, B \vdash C\)}
%     \RL{\(\lor E\) }
%     \TIC{\(\G \vdash C\)}
% \end{prooftree}


\end{itemize}
\section{Linearna logika}
    \begin{itemize}
        \item kako bi modelirali kemijski primer
        \begin{equation*}
            2 \mathsf{H_2} + \mathsf{O_2} \rightarrow 2 \mathsf{H_2 O}
        \end{equation*}
        
        \begin{itemize}
            \item Očitno moklekule niso izjave
            \item Viri in cilj
        \end{itemize}
    \end{itemize}



\subsection{Nekaj operatorjev}
        \begin{itemize}

            \item hkratna konjunkcija
            
            \begin{itemize}
                \item Različni konteksti
            \end{itemize}
\begin{prooftree}
    \AXC{\(\G \vdash A\)}
    \AXC{\(\D \vdash B\)}
    \RL{\(\otimes\)I}
    \BIC{\(\G, \D \vdash A \otimes B\)}
\end{prooftree}

\begin{prooftree}
    \AXC{\(\G \vdash A \otimes B\)}
    \AXC{\(\D, A, B \vdash C\)}
    \RL{\(\otimes E\)}
    \BIC{\(\G, \D \vdash C\)}
\end{prooftree}


\begin{prooftree}
    \AXC{2€ \rampa bela kava}
    \AXC{1,40€ \rampa piškot}
    \BIC{3,40€ \rampa bela kava \(\otimes\) piškot}
\end{prooftree}

\item Seveda!
\begin{prooftree}
    \AXC{!Mafija, 2€ \rampa bela kava}
    \AXC{!Mafija, 1,40€ \rampa piškot}
    \BIC{!Mafija, 3,40€ \rampa bela kava \(\otimes\) piškot}
\end{prooftree}

    \begin{itemize}
        \item Delinearizira
        \item ! --- Imamo dovolj stavri
        \item ? --- lahko sprejmemo dovolj stvari
    \end{itemize}



            \item linearna implikacija
\begin{prooftree}
    \AXC{\(\G, A \vdash B\)}
    \RL{\(\multimap I\)}
    \UIC{\(\G \vdash A \multimap B\)}
\end{prooftree}

\begin{prooftree}
    \AXC{\(\G \vdash A\)}
    \AXC{\(\D \vdash A \multimap B\)}
    \BIC{\(\G, \D \vdash B\)}
\end{prooftree}
%

\begin{prooftree}
    \AXC{bon \rampa sendvič}
    \UIC{\rampa bon \(\multimap\) sendvič}
\end{prooftree}




% \begin{prooftree}
%     \AXC{}
%     \AXC{2€, kartica \rampa bela kava}
%     \AXC{2€, gotovina \rampa bela kava}

% \end{prooftree}

            \item alternativna konjunkcija (notranja izbira)
\begin{prooftree}
    \AXC{\(\G \vdash A\)}
    \AXC{\(\G \vdash B\)}
    \RL{\(\& I\)}
    \BIC{\(\G \vdash A \& B\)}
\end{prooftree}
\begin{align*}
    &
    \begin{bprooftree}
        \AXC{\(\G \vdash A \& B\)}
        \RL{\(\& E_L\)}
        \UIC{\(\G \vdash A\)}
    \end{bprooftree}
    &
    \begin{bprooftree}
        \AXC{\(\G \vdash A \& B\)}
        \RL{\(\& E_D\)}
        \UIC{\(\G \vdash B\)}
    \end{bprooftree}
    &
\end{align*}


\begin{prooftree}
    \AXC{bon \rampa sendvič}
    \AXC{bon \rampa tortilija}
    \BIC{bon \rampa sendvič \& tortilija}
\end{prooftree}



           \item disjunkcija (zunanja izbira)
\begin{align*}
    \begin{bprooftree}
        \AXC{\(\G \vdash A\)}
        \RL{\(\oplus I_L\)}
        \UIC{\(\G \vdash A \oplus B\)}
    \end{bprooftree}
    \qquad
    \begin{bprooftree}
        \AXC{\(\G \vdash B\)}
        \RL{\(\oplus I_D\)}
        \UIC{\(\G \vdash A \oplus B\)}
    \end{bprooftree}
\end{align*}


\begin{prooftree}
    \AXC{\(\G \vdash A \oplus B\)}
    \AXC{\(\D, A \vdash C\)}
    \AXC{\(\D, B \vdash C\)}
    \RL{\(\oplus E\)}
    \TIC{\(\G, \D \vdash C\)}
\end{prooftree}


            \item multiplikativna disjunkcija
            
            \begin{itemize}
                \item Dual disjunkcije (negacija, multiplikativnost, aditivnost, ...)
                \item Dual ! --- ?
                \item Sekvenčni račun
            \end{itemize}

    % \subsection{Kosilo}
    %     \begin{enumerate}
    %         \item zelena solata ali rukola (\(S \& R\))
    %         \item ričet ali dunajski zrezek (\(E \& D\))
    %         \item tortica ali presenečenje šefa (palačinke ali jabolčni zavitek) (\(T \& (P \oplus J)\))
    %     \end{enumerate}
    %     \((S \& R) \otimes (E \& D) \otimes (T \& (P \oplus J))\)
    %     %
    
\end{itemize}
\pagebreak


\section{Linearni tipi}

\begin{itemize}
    \item  Curray-Howardova korespondenca (primer zapisa)
    \begin{itemize}
        \item vsi smo imeli prograiranje 1, razen profesor Strle. Sočustvujem z vami, da ste zamudili tako dobr predmet.
    \end{itemize}
    
    \begin{itemize}
        \item leva in desna stran --- Kaj pomeni zapis na eni in kaj na drugi strani
        
    \end{itemize}



    \item Primer kako dokažemo stvari hitrje v linearni logiki
\\
\\
\\
\\
\\
    \item primer z arreji (tabelami)
\\
\\
\\
\\
\\
    \item primeri programskih jezikov, ki delujejo linearno
\\
\\
\\
\\
\\
    \item rust in linearni tipi
\end{itemize}

\end{document}